%%%%%%%%%%%%%%%%%%%%%%%%%%%%%%%%
%%%%%%%%%%%%%%%%%%%%%%%%%%%%%%%%
\begin{frame}  \small \begin{ejercicio}

Muestra que el ideal $\langle 1+x \rangle$ en $V^5[x]$ correspondiente al código en $\mathbb{F}_2^5$ que tiene todas las palabras de peso par.


\end{ejercicio}
\end{frame} 

%%%%%%%%%%%%%%%%%%%%%%%%%%%%%%%%
\begin{frame}  \tiny
\[
    \begin{array}{cccl}
					 \mathbb{F}_2^3 & p(x) &  p(x)f(x) mod (x^5-1)  \\ 
					        &     &  &\\ 
00000 & 0 & \mathbf{0} 																						&00000\\
10000 & 1 & \mathbf{1+x}																					&11000\\
01000 & x & x+x^2=\mathbf{x+x^2}																	&00110\\
11000 & x+1 & 1+x+x+x^2 =\mathbf {1+x^2}													&10100 \\
00100 & x^2 & x^2+x^3 = \mathbf {x^2+x^3} 													&00110\\
10100 & x^2+1 & x^2+1+x^3+x =\mathbf {1+x+x^2+x^3}								&11110\\
01100 & x^2+x & x^2+x+x^3+x^2=\mathbf {x+x^3} 										&01010 \\
11100 & x^2+x+1 & x^2+x+1+x^3+x^2+x =\mathbf {1+x^3}						&01001 \\
& & &\\

00010 & x^3 & \mathbf{x^4+x^3}	& 																	00011\\
10010 & x^3+1 & x^3+1+x^4+x=\mathbf{1+x+x^3+x^4}& 							11011\\
01010 & x^3+x & x^3+x+x^4+x^2=\mathbf{x+x^2+x^3+x^4}&						01111\\
01110 & x^3+x+1 & x^3+x+1+x^4+x^2+x =\mathbf {x^4+x^3+x^2+1} &	10111 \\
00110 & x^3+x^2 & x^3+x^2+x^4+x^3 = \mathbf {x^4+x^2} &						00101\\
11110 & x^3+x^2+1 & x^3+x^2+1+x^4+x^3+x =\mathbf {x^4+x^2+x+1} &11101\\
01110 & x^3+x^2+x & x^3+x^2+x+x^4+x^3+x^2=\mathbf {x^4+x} 			&01001 \\
11110 & x^3+x^2+x+1 & x^3+x^2+x+1+x^4+x^3+x^2+x =\mathbf {x^4+1} &10001 \\
& & &\\

	\end{array}	    		
\]
\end{frame} 


%%%%%%%%%%%%%%%%%%%%%%%%%%%%%%%%
\begin{frame}  \tiny
\[
    \begin{array}{cccl}
					 \mathbb{F}_2^3 & p(x) &  p(x)f(x) mod (x^5-1)  \\ 
					        &     &  &\\ 
00001	 & x^4 & x^5+x^4 =\mathbf{1+x^4}																		  &10001\\
10001 & x^4+1 & x^4+1+x^5+x =\mathbf{x+x^4}															  &01001\\
01001 & x^4+x & x^4+x+x^5+x^2=\mathbf{1+x+x^2+x^4}											  &11101\\
11001 & x^4+x+1 & x^4+x+1+x^5+x^2+x =\mathbf {x^4+x^2} 									  &00101 \\
00101 & x^4+x^2 & x^4+x^2+x^5+x^3 = \mathbf {x^4+x^3+x^2+1} 							  &10111\\
10101 & x^4+x^2+1 & x^4+x^2+1+x^5+x^3+x =\mathbf {x^4+x^3+x^2+x} 				  &01111\\
01101 & x^4+x^2+x & x^4+x^2+x+x^5+x^3+x^2=\mathbf {x^4+x^3+x+1} 				  &11011 \\
10101 & x^4+x^2+x+1 & x^4+x^2+x+1+x^5+x^3+x^2+x =\mathbf {x^4+x^3}			  &11000 \\
& & &\\
00011 & x^4+x^3 & x^4+x^3+x^5+x^4 =\mathbf{1+x^3}												   &10010\\
01011 & x^4+x^3+1 & x^4+x^3+1+x^5+x^4+x =\mathbf{x+x^3}									   &01010\\
01011 & x^4+x^3+x & x^4+x^3+x+x^5+x^4+x^2=\mathbf{1+x+x^2+x^3}					   &11110\\
11011 & x^4+x^3+x+1 & x^4+x^3+x+1+x^5+x^4+x^2+x =\mathbf {x^3+x^2} 			   &00110\\
00111 & x^4+x^3+x^2 & x^4+x^3+x^2+x^5+x^4+x^3 = \mathbf {x^2+1} 					   &10100\\
10111 & x^4+x^3+x^2+1 & x^4+x^3+x^2+1+x^5+x^4+x^3+x =\mathbf {x^2+x} 		   &01100\\
01111 & x^4+x^3+x^2+x & x^4+x^3+x^2+x+x^5+x^4+x^3+x^2=\mathbf {x+1} 		   &11000\\
11111 & x^4+x^3+x^2+x+1 & x^4+x^3+x^2+x+1+x^5+x^4+x^3+x^2+x =\mathbf {0} &00000\\
& & &\\


	\end{array}	    		
\]
\end{frame}

%%%%%%%%%%%%%%%%%%%%%%%%%%%%%%%%
\begin{frame}  \small \begin{ejercicio} 

9.11. La factorización de $x^{15}-1$ en $\mathbb{F}_2[x]$ es

 \begin{center}
$x^{15}-1 = (1+x)(1+x+x^2)(1+x+x^4)(1+x^3+x^4)(1+x+x^2+x^3+x^4). $
\end{center}
 
 Utilizando este resultado, encuentra un generador canónico para un código cíclico equivalente al código Hamming de longitud 15.
 

\end{ejercicio} \pause \textit{\underline{Solución}}:\\ \end{frame} 

%%%%%%%%%%%%%%%%%%%%%%%%%%%%%%%%
%%%%%%%%%%%%%%%%%%%%%%%%%%%%%%%%
%SECCION %%%%%%%%%%%%%%%%%%%%%%%%%%%%%%%
%%%%%%%%%%%%%%%%%%%%%%%%%%%%%%%%
%%%%%%%%%%%%%%%%%%%%%% 
\section{Códigos que corrigen más que un error}

%%%%%%%%%%%%%%%%%%%%%%%%%%%%%%%%
\begin{frame}  \frametitle{Códigos que corrigen más que un error}\small  

La siguiente generalización del teorema \ref{teo:Hdef} es el punto de comienzo para la construcción de familias de códigos con una distancia mínima arbitrariamente grande.

\begin{teorema}

Sea H una matriz de verificación para un código binario C. La mínima distancia de C es igual al mínimo número de columnas linealmente dependiente de H.

\end{teorema}


\end{frame} 


%%%%%%%%%%%%%%%%%%%%%%%%%%%%%%%%
\begin{frame}  \small  

\begin{itemize} 
\item  El teorema indica que una matriz de verificación es el cual cada conjunto de r columnas es linealmente independiente define un código con una distancia mínima de al menos r+1. 

\item El siguiente lema muestra que tipo particular de matriz tiene esta propiedad, y sea usado para justificar la construcción que veremos en la próxima sección.
\end{itemize}


\begin{lema}

F es un cuerpo y $a \in F$   es tal que $a^0=1,a,a^2,a^3,\ldots , a^{n-1}$ son elementos distintos y no ceros de F. \\
\vspace*{5pt}
Sea A la matriz r x n tal que la entrada en fila i y columna j es $a^{ij}$, donde las filas son etiquetadas $ i = 1,2,\ldots ,r$ y las columnas son etiquetadas $ j = 0,1,\ldots ,n-1$. \\
\vspace*{5pt}
Entonces cualquier r columnas de A son linealmente independiente sobre F.\\

\end{lema}
\end{frame} 
%%%%%%%%%%%%%%%%%%%%%%%%%%%%%%%%


%%%%%%%%%%%%%%%%%%%%%%%%%%%%%%%%
%%%%%%%%%%%%%%%%%%%%%%%%%%%%%%%%
%\begin{frame}  \small \begin{ejercicio} \end{ejercicio} \pause \textit{\underline{Solución}}:\\ \end{frame} 
%%%%%%%%%%%%%%%%%%%%%%%%%%%%%%%%
%%%%%%%%%%%%%%%%%%%%%%%%%%%%%%%%
%EJERCICIO %%%%%%%%%%%%%%%%%%%%%%%%%%%%%%%
%%%%%%%%%%%%%%%%%%%%%%%%%%%%%%%%

%%%%%%%%%%%%%%%%%%%%%%%%%%%%%%%%
\begin{frame}  \small \begin{ejercicio} 

9.15. Supongase que queremos asignar números ID de la forma de palabras binarias de longitud n a un conjunto de 16 personas, por lo que las palabras forman un  código lineal corrector 2, que es un código lineal con parámetros (n,4,5). Muestra que necesitamos que n sea al menos 10.

\end{ejercicio} \pause \textit{\underline{Solución}}:\\ \end{frame} 


%9.16. Given the conditions stated in the previous exercise, there is in fact no solution with n = 10. [You are not asked to prove this.] However, there is a solution with n = 11, which we are going to construct. Let A be the linear code with parameters (6, 3, 3) discussed in Example 8.6 and Section 8.3), and let B = {000000, 111111} be the code with parameters (6, 1, 6).  Let $C \subseteq \mathbb{F}_2^{12}$ be the code formed by words of the form [a | a + b] $a \in A$, $b \in B$. Show that C is a linear code with parameters (12, 4, 6) and construct a linear code with parameters (11, 4, 5) by making a suitable modification to C.


%9.17. Is the code constructed in the previous exercise a cyclic code?




%%%%%%%%%%%%%%%%%%%%%%%%%%%%%%%%
%\begin{frame}  \small  
%\begin{itemize} \item  \item \item \item 
%\end{itemize} \end{frame} 


%%%%%%%%%%%%%%%%%%%%%%%%%%%%%%%%
%\begin{frame}  \small  
%\end{frame} 
%%%%%%%%%%%%%%%%%%%%%%%%%%%%%%%%


%%%%%%%%%%%%%%%%%%%%%%%%%%%%%%%%
%%%%%%%%%%%%%%%%%%%%%%%%%%%%%%%%
%SECCION %%%%%%%%%%%%%%%%%%%%%%%%%%%%%%%
%%%%%%%%%%%%%%%%%%%%%%%%%%%%%%%%
%%%%%%%%%%%%%%%%%%%%%% 
\section{Definición de una familia de códigos BCH}

%%%%%%%%%%%%%%%%%%%%%%%%%%%%%%%%
\begin{frame}  {Definición de una familia de códigos BCH}\small  

\begin{itemize} 

\item  Las familias de códigos que pueden corregir cualquier número de errores fueron descritos por Bose y Ray-Chaudhuri, e (independientemente) por Hocqhenghem, en 1959-60. Estos códigos son conocidos como códigos BCH.

\item Aquí solo consideramos un solo caso: una familia de códigos cíclicos binarios de tal forma que, dado los enteros m y t, los parámetros $(n,k,\delta)$ satisface:

\begin{center}
$n=2^m-1$,	$k\geq n-tm$, 	$\delta\geq 2t+1.$
\end{center}

\item Ya que los BCH códigos son códigos cíclicos, comenzamos con la factorización de $x^n -1$ en polinomios irreducibles. Supongamos que 

\begin{center}
$x^n - 1 = (x - 1)f_1(x)f_2(x) \ldots  f_l(x)	en \mathbb{F}_2[x]$.
\end{center}

\end{itemize} 

\end{frame} 

%%%%%%%%%%%%%%%%%%%%%%%%%%%%%%%%
\begin{frame}  \small  
\begin{itemize} 

\item  Hay solamente un posible factor lineal x-1, y cuando n es impar ocurre solamente una vez, por lo que los polinomios $f_1(x), \ldots , f_l(x)$ tienen todos un grado mayor que 1. 

\item Sin embargo, si extendemos el cuerpo entonces $x^n-1$ puede ser dividido en factores lineales, o en otras palabras, tiene n raíces. 

\item Nuestra definición de los códigos BCH supone la construcción de un cuerpo E que contiene $\mathbb{F}_2$ y un elemento $\alpha in E$ tal que $\alpha$ es una raíz de las ecuaciones $x^n=1$. 

\item Predisponemos que los elementos de E son distintos, en el que ese caso diremos que $\alpha$ una raíz primitiva. 

\item La construcción de $\alpha$ y E la estableceremos en breve, pero por ahora nos basaremos en el supuesto que puede ser realizado. Esto significa que podemos escribir:

\begin{center}
$x^n -1 = (x-1)(x-\alpha)(x-\alpha^2) \ldots (x-\alpha^{n-1})$ en $E[x]$.
\end{center}

\end{itemize} \end{frame} 


%%%%%%%%%%%%%%%%%%%%%%%%%%%%%%%%
\begin{frame}  \small  

\begin{lema}
Si $\alpha$ es una raíz primitiva, cada uno de los términos $x-\alpha^{i} (1\leq i \leq n-1)$ es un factor de exactamente uno de los polinomios $f_j(x) (j = 1,2,\ldots ,l)$ en $E[x]$.
\end{lema}

Para $i = 1,2,\ldots ,n-1$ sea $m_i(x)$ el único factor irreducible $f_j(x)$ de $x^n-1$ en $\mathbb{F}_2[x]$ para el cual $x-\alpha^i$ es un factor de $f_j(x)$ en E[x]. Equivalentemente, $m_i(x)=f_j(x)$ cuando $f_j(x)$ tiene $\alpha^i$ como raíz.

\begin{definicion}[Código BCH, Distancia Diseñada]
Dado un entero impar n, sea los polinomios $m_i(x) \in \mathbb{F}_2[x]$ definido anteriormente, y definimos g(x) como el menor múltiplo común de 


\begin{center}
$m_1(x), m_2(x), \ldots ,m_{dd-1}(x). $
\end{center}

Decimos que el código cíclico en $\mathbb{F}_2^n$ con un generador canónico g(x) es un código BCH con la distancia diseñada $dd$.
\end{definicion}

\end{frame} 
%%%%%%%%%%%%%%%%%%%%%%%%%%%%%%%%


%%%%%%%%%%%%%%%%%%%%%%%%%%%%%%%%
\begin{frame}  \small  
\begin{itemize} 

\item  La razón para tomar el mínimo común múltiplo (m.c.m.) es que los polinomios $m_i(x)$ puede no ser todos distintos.

\item Por ejemplo, para cualquier $f(x) \in \mathbb{F}_2[x]$ tenemos que $f(x)^2=f(x^2)$. 

\item Seguimos que si $\alpha^i$ es una raíz de $f_j(x)$ entonces tenemos $\alpha^{2i}$.

\item Por tanto los polinomios $m_i(x)$ y $m_{2i}(x)$ son los mismo, y en la definición de g(x) consideramos solamente los polinomios $m_i(x)$ con i impar. 

\item Cuando dd es un número impar, la definición de g(x) es:

\begin{center}
$g(x) = mcm 	\{m_1(x), m_3(x), \ldots , m_{dd-2}(x)\}.$
\end{center}

\item Abordamos el problema de construir la raíz primitiva $\alpha$ y el cuerpo E. Considera el caso n=7, cuando los factores irreducibles de $x^7-1$ son

\begin{center}
$x^7 -1=(1+x)(1+x+x^3)(1+x^2 +x^3).$
\end{center}

\end{itemize} 


\end{frame} 

%%%%%%%%%%%%%%%%%%%%%%%%%%%%%%%%
\begin{frame}  \footnotesize

\begin{itemize} 


\item  Cualquier raíz de $x^7-1$ debe tener una raíz de uno de los factores. 

\item La única raíz de 1+x, es 1, el cual no es primitiva. 

\item Si $\alpha$ es una raíz de $1+x+x^3$, la ecuación $1+\alpha+\alpha^3=0$ puede ser escrita como $\alpha^3=1+\alpha$. 

\item Utilizando esta ecuación, cualquier potencia de $\alpha$ puede ser reducida a un polinomio en $\alpha$ con grado no mayor que 2. Por tanto:
\[
     \begin{array}{l}
     				\alpha^3	=1+\alpha\\
					\alpha^4	=\alpha (1+\alpha)=\alpha +\alpha^2\\
					\alpha^5	=\alpha (\alpha +\alpha^2)=\alpha^2 +\alpha^3 =1+\alpha +\alpha^2\\
					\alpha^6	=\alpha (1+\alpha +\alpha^2)=1+\alpha^2\\
					\alpha^7	=\alpha(1+\alpha^2)=1.
           			\end{array} 		
\]
\end{itemize} 



Utilizando la correspondencia: \\

\begin{center}
$c_0+c_1\alpha +c_2\alpha^2 \leftrightarrow c_0c_1c_2$
\end{center}

tenemos\\

\begin{center}
$\alpha	= 010,	\alpha^2	= 001,	\alpha^3	= 110,	\alpha^4	= 011 $\\
$\alpha^5	= 111,	\alpha^6	= 101,	\alpha^7	= 100.$
\end{center}

\end{frame} 
%%%%%%%%%%%%%%%%%%%%%%%%%%%%%%%%
\begin{frame}  \footnotesize 
\begin{itemize} 


\item  Hemos representado las 7 potencias de $\alpha$ como los 7 elementos distintos y no ceros de $\mathbb{F}_2^3$. Esto significa que podemos tomar E como $\mathbb{F}_2^3$, con las apropiadas operaciones de multiplicación: el producto de $a_0a_1a_2$ y $b_0b_1b_2$ es obtenido multiplicando los polinomios $a_0 + a_1 \alpha+ a_2 \alpha^2$ y  $b_0 + b_1 \alpha+ b_2 \alpha^2$ y utilizando la relación $1 + \alpha + \alpha^3 = 0$ para reducir al polinomio.

\item El mismo método funciona con cualquier $n=2^m -1$: es posible encontrar uno de los factores irreducibles de $x^n-1$ que define una raíz primitiva $\alpha$.  %The reader who has studied Galois fields will recognise the construction of GF(2m).

\item Para resumir: cuando $n=2^m-1$ el método descrito anteriormente puede ser siempre utilizado para construir un cuerpo E con $2^m$ elementos, y una raíz primitiva $\alpha \in E$. 

\item Además, como veremos posteriormente, el mismo método produce una matriz de verificación para un código BCH con $n=2^m-1$ y un valor dado de $dd$.

\item Mostraremos un ejemplo, donde la matriz de verificación puede ser obtenida directamente.


\end{itemize} \end{frame} 


%%%%%%%%%%%%%%%%%%%%%%%%%%%%%%%%
\begin{frame}  \small  

\begin{ejemplo}

Tomaremos n = 7 y $\alpha$ una raíz de $1+x+x^3$. Encuentre el generador canónico para el código BCH de una distancia diseñada dd=7, y muestra que el código tiene una distancia mínima $\delta=7$.\\
\end{ejemplo}

\end{frame} 
%%%%%%%%%%%%%%%%%%%%%%%%%%%%%%%%

%%%%%%%%%%%%%%%%%%%%%%%%%%%%%%%%
\begin{frame}  \small  

\textit{\underline{Solución}}:\\

\begin{itemize} 


\item  Ya que $dd = 7$, el generador canónico es el menor común múltiplo de $m_1(x)$,$m_3(x)$ y $m_5(x)$. Tenemos que $\alpha$ es una raíz de $1+x+x^3$, por lo que $m_1(x) = 1+x+x^3$.

\item Para determinar $m_3(x)$ y $m_5(x)$ necesitamos encontrar los factores que tiene $\alpha^3$ y $\alpha^5$ como raíces. Utilizando la representación como elementos de $\mathbb{F}_2^3$ tenemos:

\begin{center}
$1+(\alpha^3)^2 +(\alpha^3)^3 =1+\alpha^6 +\alpha^2 =100+101+001=000,$\\
$1+(\alpha^5)^2 +(\alpha^5)^3 =1+\alpha^3 +\alpha=100+110+010=000.$
\end{center}


\end{itemize} 

\end{frame} 
%%%%%%%%%%%%%%%%%%%%%%%%%%%%%%%%

%%%%%%%%%%%%%%%%%%%%%%%%%%%%%%%%
\begin{frame}  \small  

\textit{\underline{cont Solución}}:\\

\begin{itemize} 

\item En otras palabras, ambos $\alpha^3$ y $\alpha^5$ son raíces de $1+x^2+x^3$ y $m_3(x)=m_5(x)=1+x^2+x^3$. Por tanto,

\begin{center}
$g(x) = mcm \{m_1(x), m_3(x), m_5(x)\} = (1 + x + x^3)(1 + x^2 + x^3).$
\end{center}

 \item Siguiendo las reglas dadas anteriormente, encontramos h(x) = 1 + x, por lo que la dimensión es k=1. La primera fila de la matriz de verificación es $h^* = 1100000$ y las otras filas son los desplazamientos cíclicos $0110000, \ldots, 0000011$. 
 
 \item Las ecuaciones resultantes implican que $x_1x_2\ldots x_7$ es una palabra codificada sii $x_1=x_2=\ldots x_7$, por lo que el código es \{0000000, 1111111\}. Claramente, este código tiene una distancia mínima de 7.

\end{itemize} 

\end{frame} 
%%%%%%%%%%%%%%%%%%%%%%%%%%%%%%%%




%%%%%%%%%%%%%%%%%%%%%%%%%%%%%%%%
%\begin{frame}  \small  
%\begin{itemize} \item  \item \item \item 
%\end{itemize} \end{frame} 


%%%%%%%%%%%%%%%%%%%%%%%%%%%%%%%%
%\begin{frame}  \small  
%\end{frame} 
%%%%%%%%%%%%%%%%%%%%%%%%%%%%%%%%


%%%%%%%%%%%%%%%%%%%%%%%%%%%%%%%%
%%%%%%%%%%%%%%%%%%%%%%%%%%%%%%%%
%\begin{frame}  \small \begin{ejercicio} \end{ejercicio} \pause \textit{\underline{Solución}}:\\ \end{frame} 
%%%%%%%%%%%%%%%%%%%%%%%%%%%%%%%%
%%%%%%%%%%%%%%%%%%%%%%%%%%%%%%%%
%EJERCICIO %%%%%%%%%%%%%%%%%%%%%%%%%%%%%%%
%%%%%%%%%%%%%%%%%%%%%%%%%%%%%%%%

%9.18. Prove that f(x)2 = f(x2) for any polynomial f(x) # F2[x]. Hence
%determine all the polynomials mi(x), 1 - i - 6 when n = 7 and the
%primitive root is a root of 1 + x + x3.

%%%%%%%%%%%%%%%%%%%%%%%%%%%%%%%%
%%%%%%%%%%%%%%%%%%%%%%%%%%%%%%%%
\begin{frame}  \small \begin{ejercicio} 

9.19. Factoriza $x^3-1$ en $\mathbb{F}_2[x]$, y por tanto determina el código binario BCH en $\mathbb{F}_2^3$ con dd=3.  (Refer to Exercise 9.5.)

\end{ejercicio} \pause \textit{\underline{Solución}}:\\ \end{frame} 


%%%%%%%%%%%%%%%%%%%%%%%%%%%%%%%%
%%%%%%%%%%%%%%%%%%%%%%%%%%%%%%%%
\begin{frame}  \small \begin{ejercicio} 

9.20. Muestra que la raíz de  $1 + x^2 + x^3$ puede ser tomada como la raíz primitiva de $x^7 - 1$.

\end{ejercicio} \pause \textit{\underline{Solución}}:\\ \end{frame} 




%%%%%%%%%%%%%%%%%%%%%%%%%%%%%%%%
%%%%%%%%%%%%%%%%%%%%%%%%%%%%%%%%
%SECCION %%%%%%%%%%%%%%%%%%%%%%%%%%%%%%%
%%%%%%%%%%%%%%%%%%%%%%%%%%%%%%%%
%%%%%%%%%%%%%%%%%%%%%% 
\subsection{Propiedades de los códigos BCH}

%%%%%%%%%%%%%%%%%%%%%%%%%%%%%%%%
\begin{frame}  {Propiedades de los códigos BCH}\small  

\begin{itemize} 
\item  Supongamos que conocemos una raíz primitica $\alpha$ de $x^n-1$. 

\item El siguiente lema nos permite construir una matriz de verificación para el correspondiente código BCH sin tener que calcular el generador canónico explicitamente. 
\end{itemize} 


\begin{lema}

Sea $n = 2^m -1$. Si $c = c_0c_1\ldots c_{n-1}$ es una palabra codificada en un código binario BCH de distancia diseñada dd=2t+1, construido utilizando una raíz primitiva $\alpha$, entonces $H^bc'=0'$, donde $H^b$ es la matriz 2t x n sobre E dada por 

\[
    H^b =   \left( \begin{array}{ccccccc}
     				  1 &  \alpha  & \alpha^2  & \ldots & \alpha^{n-1}\\ 
					  1 &  \alpha^2  & \alpha^4  & \ldots & \alpha^{2(n-1)}\\ 
					  . & .   & .  & \ldots & .\\ 
					  1 &  \alpha^{2t}  & \alpha^{4t}  & \ldots & \alpha^{2t(n-1)}
           			\end{array}
    		\right) 	    		
\]

\end{lema}
\end{frame} 

%%%%%%%%%%%%%%%%%%%%%%%%%%%%%%%%
\begin{frame}  \small  

\begin{teorema}\label{te:dm}

Sea C un código binario BCH con longitud de palabra $n=2^m -1$ y la distancia diseñada $dd = 2t + 1$. Entonces la dimensión k y la \textit{distancia mínima} $\delta$ de C satisface 

\begin{center}
$k \geq n - tm$,	$\delta \geq dd$.
\end{center}

\end{teorema}
\end{frame} 


%%%%%%%%%%%%%%%%%%%%%%%%%%%%%%%%
\begin{frame}  \small  

\begin{ejemplo}

La factorización  de $x^{15} -1$ en $\mathbb{F}_2[x]$ es:

$x^{15} -1=(1+x)(1+x+x^2)(1+x+x^4)(1+x^3 +x^4)(1+x+x^2 +x^3 +x^4).$

Muestra que una raíz $\alpha$ de $1 + x + x^4$ puede ser escogida como raíz primitiva, y por tanto desarrolla una matriz de verificación para el código BCH con $n = 15$ y $dd = 5$.\\
\end{ejemplo}

\end{frame} 

%%%%%%%%%%%%%%%%%%%%%%%%%%%%%%%%
\begin{frame}  \small  

\textit{\underline{Solución}}:\\

Ya que $\alpha^4 = 1 + \alpha$, cualquier potencia de $\alpha$ puede ser reducido a un polinomio en $\alpha$ con un grado de al menos 3. Utilizando la correspondencia 

\begin{center}
$c_0 +c_1\alpha +c_2 \alpha^2 +c_3\alpha^3 \leftrightarrow c_0 c_1 c_2 c_3$
\end{center}

encontramos que:\\


\begin{center}
$\alpha = 0100, \alpha^2 = 0010,	\alpha^3 = 0001, \alpha^4 = 1100,$\\
$\alpha^5 = 0110, \alpha^6 = 0011, \alpha^7 = 1101, \alpha^8 = 1010,$\\
$\alpha^9 = 0101, \alpha^{10} = 1110, \alpha^{11} = 0111, \alpha^{12} = 1111,$\\
$\alpha^{13} = 1011, \alpha^{14} = 1001, \alpha^{15} = 1000.$\\

\end{center}


\end{frame} 

%%%%%%%%%%%%%%%%%%%%%%%%%%%%%%%%
\begin{frame}  \small  

\textit{\underline{continua Solución}}:\\


Ya que todos los vectores son distintos, concluimos que $\alpha$ es una raíz primitiva. En el caso $dd=5$, g(x) es el mcm de $m_1(x),m_2(x),m_3(x).$ Ya que $m_2(x) =m_1(x)$ una matriz de verificación para el código es la matriz 8 x 15 sobre $\mathbb{F}_2$, 

\[
    H=   \left( \begin{array}{ccccccccc}
     				  1 &  \alpha  & \alpha^2  & \alpha^3 & \alpha^4 & \alpha^5 & \alpha^6  &\ldots & \alpha^{14}\\ 
					  1 &  \alpha^3  & \alpha^6  & \alpha^9 & \alpha^{12} & 1& \alpha^3  &\ldots & \alpha^{12}\\ 
          			\end{array}
    		\right) 	    		
\]

donde $\alpha^ i$ representa el correspondiente vector columna, listado anteriormente.\\
\vspace*{5pt}
Según el teorema \ref{te:dm}, el código construido anteriormente tiene $\delta \geq dd =5$ y por tanto es un código corrector de error 2. Existe un procedimiento simple para corrección de errores:


\end{frame} 
%%%%%%%%%%%%%%%%%%%%%%%%%%%%%%%%
\begin{frame}  \small  
\textit{\underline{continua Solución}}:\\


Suponemos que recibimos una palabra $z= z_0z_1 \ldots z_{14}$. Como es usual calculamos el síndrome s,  $s'= Hz' = [x y]'$, donde x e y son vectores de tamaño 4.

\begin{itemize}
\item Si x=y=0000 entonces z es una palabra codificada y no hay errores en los bits.

\item Si, para algún i, $x=\alpha^i$ e $y=\alpha^{3i}$ el síndrome es la $i^{th}$ columna de H, y existe un error en el $i^{th}$ bit de z.

\item Si es posible encontrar i y j que cumplan las ecuaciones 

\begin{center}
$x=\alpha^i +\alpha^j$, 	$y=\alpha^{3i} +\alpha^{3j},$
\end{center}

entonces el síndrome es la suma de las columnas $i^{th}$ y $j^{th}$ de H, y z tiene errores en los bits $i^{th}$ y $j^{th}$. Por tanto debemos resolver las ecuaciones para obtener i y j. Si las ecuaciones no tienen solución, debemos asumir que más de errores se han producido.
\end{itemize}

\end{frame} 


%%%%%%%%%%%%%%%%%%%%%%%%%%%%%%%%
\begin{frame}  \small  

\begin{ejemplo}

Supongamos que dos errores se produce en la transmisión y el síndromde la palabra recibida z es 

\begin{center}
$s' = Hz' = 1100 0110'.$
\end{center}

¿En que bit se ha producido el error?\\
\end{ejemplo}
\end{frame} 


%%%%%%%%%%%%%%%%%%%%%%%%%%%%%%%%
\begin{frame}  \small  
\textit{\underline{Solución}}:\\

En este caso, x = 1100 e y = 0110. Si los errores han sido producidos en los bits $i^{th}$ y $j^{th}$, tenemos que encontrar i y j tal que 

\begin{center}
$\alpha^i + \alpha^j = 1100$	  $\alpha^{3i} + \alpha^{3j} = 0110.$
\end{center}

Una forma de resolverlo es realizar una lista de pares $(\alpha^i, \alpha ^j)$, para los cuales la primera ecuación se cumple:

\begin{center}
$(1,\alpha), (\alpha,1), (\alpha^2,\alpha^{10}) (\alpha^3,\alpha^7), \ldots$ .
\end{center}

\end{frame} 

%%%%%%%%%%%%%%%%%%%%%%%%%%%%%%%%
\begin{frame}  \small  
\textit{\underline{continua Solución}}:\\

Ahora podemos verificar si el correspondiente par cumple la segunda ecuación:

\begin{center}
$1+\alpha^3 =1000+0001=1001,$\\
$\alpha^3 +1=0001+1000=1001,$\\
$\alpha^6 +1=0011+1000=1011,$\\
$\alpha^9 +\alpha^6 =0101+0011=0110.$\\
\end{center}

La conclusión es que los bits 3 y 7 son erróneos, por lo que la palabra codificada era

\begin{center}
$c = z + 000100010000000.$
\end{center}


\end{frame} 



%%%%%%%%%%%%%%%%%%%%%%%%%%%%%%%%
%%%%%%%%%%%%%%%%%%%%%%%%%%%%%%%%
%\begin{frame}  \small \begin{ejercicio} \end{ejercicio} \pause \textit{\underline{Solución}}:\\ \end{frame} 
%%%%%%%%%%%%%%%%%%%%%%%%%%%%%%%%
%%%%%%%%%%%%%%%%%%%%%%%%%%%%%%%%
%EJERCICIO %%%%%%%%%%%%%%%%%%%%%%%%%%%%%%%
%%%%%%%%%%%%%%%%%%%%%%%%%%%%%%%%


%%%%%%%%%%%%%%%%%%%%%%%%%%%%%%%%
%%%%%%%%%%%%%%%%%%%%%%%%%%%%%%%%
\begin{frame}  \small \begin{ejercicio} 

9.21. Compara la tasa de información y las propiedades correctora de erro del código BCH construido en el Ejemplo 9.22 con el código Hamming de la misma longitud de palabra.

\end{ejercicio} \pause \textit{\underline{Solución}}:\\ \end{frame} 


%%%%%%%%%%%%%%%%%%%%%%%%%%%%%%%%
%%%%%%%%%%%%%%%%%%%%%%%%%%%%%%%%
\begin{frame}  \small \begin{ejercicio} 

9.22. Supongamos que el código BCH en ele Ejemplo 9.22 es utilizado y una palabra de síndrome 10101001 es recibida. Asumiendo que ha tiene  dos bits erróneos, ¿Cuales son los bits que deben ser corregidos?

\end{ejercicio} \pause \textit{\underline{Solución}}:\\ \end{frame} 

%%%%%%%%%%%%%%%%%%%%%%%%%%%%%%%%
%%%%%%%%%%%%%%%%%%%%%%%%%%%%%%%%
\begin{frame}  \small \begin{ejercicio} 

9.23. Investiga los códigos binarios BCH con n=15 y dd=7, 9, 11, 13, 15, tomando una raíz de $1 + x + x^4$ como raíz primitiva como en el Ejemplo  9.22.

\end{ejercicio} \pause \textit{\underline{Solución}}:\\ \end{frame} 

